% AI-written, AI cross-reviewed; no human sign-off; not a claim that RH is proved
% Atlas unit D154. Local manuscript v3, item G.164, heading on printed PDF p. 198.
% ChowFormCompatibility.chow_form_hyperplane_section
% paper2/decomposition_tables.tex:1401–1406. Verbatim excerpt; SHA-256 2cce46319a3d72acec098cf7f5482b7f4223530250c143610565529a09ca745a.
% Original TeX uses the macros and cross-reference keys of the source manuscript.

\begin{formalstatement}[D154]\label{decomp:D154}
Slicing the first block of the Chow form by a hyperplane gives, up to a nonzero scalar, the product of the component Chow forms raised to the intersection multiplicities.
\end{formalstatement}
{\small\raggedright\noindent\textsf{\bfseries Related reference}\quad \nolinkurl{ChowFormCompatibility.chow_form_hyperplane_section}\par\noindent \nolinkurl{Development/WeilBound/CoreStatements/NefModelVolumes/ChowFormHyperplaneSection.lean:12--38}\par\noindent SHA-256 \nolinkurl{2a4f525036952efa0c225601133cc7b293c39ca1e923a180908e9f1e0bfa1965}\par\noindent Dominated 0\quad Disjoint 0\quad terminal.\par}
The project’s fixed-block specialization and exact intersection multiplicities require identification.
\hyperref[lit:S036]{Published sources and search S036}.
